\documentclass[12pt]{article}
\usepackage[utf8]{inputenc}
\usepackage[T1]{fontenc}
\usepackage{hyperref}
\title{Uruangnirin Language Resources for K-12 Curriculum Development}
\author{Compiled by Claude Sonnet 5 (claude-sonnet-5) for Andr\'as Kornai}
\date{}
\begin{document}
\maketitle
\section{Introduction}
Uruangnirin (ISO 639-3 urn, Glottocode urua1244) is an endangered Austronesian language of the Tanimbar-Bomberai subgroup. It is spoken by about 400 people on the Karas Islands (Tarak and Faor) off the Bomberai Peninsula, West Papua, Indonesia. Children born after about 2000 are reportedly not acquiring it.
\section{Primary Educational Resources}
None found.
\section{Linguistic and Reference Materials}
\begin{itemize}
\item Uruangnirin language (Wikipedia): \url{https://en.wikipedia.org/wiki/Uruangnirin_language}
\item E. Visser, A Phonological Sketch of Uruangnirin, an Endangered Austronesian Language of Eastern Indonesia (JSEALS): \url{https://evols.library.manoa.hawaii.edu/items/8818a0c2-240a-4550-9c48-adef0b11549b}
\item Related article in Language Documentation \& Description: \url{https://www.lddjournal.org/article/id/993}
\item Endangered Languages Project entry: \url{https://endangeredlanguages.com/lang/2798}
\end{itemize}
\section{Religious and Community Resources}
None found.
\section{Digital and Audio Resources}
The search summary mentions an online audiovisual corpus, but no direct URL for it appeared in the results, so none is cited.
\section{Curriculum Development Context}
Documentation is scholarly and recent. There is no known school use. Revitalisation work could build on the phonological sketch and the recordings.
\end{document}